% Conversão fiel da participação como cliente Hoortech; origem imutável RELATORIO_ATUAL.odt.

% SOURCE text-h-0038
\refstepcounter{section}
\section*{6.3 Conclusão}
\label{cliente-hoortech:text-h-0038}
\addcontentsline{toc}{section}{6.3 Conclusão}

% SOURCE text-p-0751
\originaltab{}Participar da AGES como cliente foi uma experiência singular e transformadora. Estar do outro lado da mesa, assumindo o papel de stakeholder, exigiu uma mudança profunda na minha mentalidade e abordagem. Foi necessário ajustar meu comportamento, minha forma de pensar e até mesmo como me comunicava com os alunos. Aprendi a dosar expectativas, a dialogar de maneira construtiva e a cobrar melhorias e novas funcionalidades com clareza e objetividade. Essa experiência me desafiou a ser um cliente que, ao mesmo tempo, inspirasse confiança e orientasse o grupo com foco no sucesso do projeto.\par

% SOURCE text-p-0752
\originaltab{}No início da jornada, enfrentamos alguns conflitos em relação ao termo de abertura do projeto, que inicialmente não refletia completamente o que queríamos alcançar. Essa divergência trouxe desafios na definição das expectativas e na comunicação das nossas demandas. No entanto, conseguimos superar esses problemas ao longo do processo, alinhando nossas visões e garantindo que todos estivéssemos no mesmo plano.\par

% SOURCE text-p-0753
Ter um projeto idealizado em apenas 24 horas, durante uma Hackatona, ampliado e desenvolvido por uma equipe altamente capacitada foi, sem dúvidas, um privilégio. A sensação de ver a ideia ganhar forma e, finalmente, torná-la um produto funcional e bem-acabado é indescritível. Mais do que apenas um resultado técnico, o que mais me marcou foi perceber o esforço, a dedicação e o talento de cada aluno refletidos em cada detalhe do trabalho entregue.\par

% SOURCE text-p-0754
Entretanto, ser cliente de um grupo de 13 alunos também trouxe grandes responsabilidades. Eles confiavam no meu suporte para esclarecer dúvidas, resolver impasses e tomar decisões importantes relacionadas às regras de negócio. Essa confiança me fez refletir profundamente sobre comprometimento e gestão, além de me ensinar lições valiosas sobre como liderar e orientar equipes em projetos de alta complexidade. Essa responsabilidade foi, sem dúvida, um fator que contribuiu para o meu amadurecimento pessoal e profissional.\par

% SOURCE text-p-0755
A oportunidade de ver o meu projeto selecionado pela AGES, entre tantas ideias, foi motivo de grande alegria para mim e para o meu time da Hackatona. Saber que nossa proposta seria levada adiante por uma equipe tão dedicada trouxe um sentimento de realização e gratidão. A confiança da AGES nos alunos e no potencial da ideia foi um grande incentivo para mim, e fico profundamente agradecido por ter tido essa chance de participar ativamente de um processo tão enriquecedor.\par

% SOURCE text-p-0756
Levo dessa experiência aprendizados que ultrapassam a esfera técnica e que, sem dúvidas, me acompanharão ao longo da vida. A vivência como stakeholder me ensinou sobre colaboração, comunicação, resiliência e, acima de tudo, sobre o valor do trabalho em equipe e do esforço conjunto. Sou imensamente grato à AGES e a todos os envolvidos por essa jornada transformadora.\par

% FIGURE draw-image-0008; source text-p-0759; image provenance pending human identification.
\memorialfigure{imagens/original/asset-0013.jpg}{\textbf{Figura 34 –} Projeto Destaque 2024/2 - CLIENTE AGES}{cliente-hoortech:draw-image-0008}{Fonte: pendente de identificação}{0.92}

